\documentclass[10pt,a4paper]{amsart}
\usepackage[margin=19mm]{geometry}
\usepackage{amsmath,amssymb,mathtools}
\usepackage[scr=rsfs,cal=euler]{mathalfa}
\usepackage{xfrac}
\usepackage[unicode,hidelinks,bookmarks=false]{hyperref}
\newcommand{\ie}{i.e.\ }
\newcommand{\cf}{cf.\ }
\newcommand{\resp}{respectively\ }
\newcommand{\Subsection}{\subsection}
\providecommand{\nobreakdash}{\nobreak\hbox{-}}
\setlength{\emergencystretch}{2em}
\multlinegap0pt
\usepackage{bm}
\usepackage{bbm}
\usepackage{dsfont}
\usepackage{xspace}
\usepackage{cancel}
\usepackage{mathtools}
\usepackage{amsthm}
\makeatletter
\@ifundefined{theorem}{\newtheorem{theorem}{Theorem}}{}
\@ifundefined{lemma}{\newtheorem{lemma}[theorem]{Lemma}}{}
\@ifundefined{proposition}{\newtheorem{proposition}[theorem]{Proposition}}{}
\makeatother
\DeclareMathOperator{\conv}{conv}
\newcommand{\RR}{\mathbb{R}}
\newcommand{\cox}{\mathrm{cox}}
\title[Toric extensions of P\'olya's theorem]{Toric extensions of P\'olya's theorem}
\author{Lorenzo Baldi, Rainer Sinn, M\'at\'e L. Telek, Julian Weigert}
\date{Source: arXiv:2511.07342v1 (CC BY 4.0)}
\begin{document}
\maketitle

\noindent\emph{Source and licence.} Toric extensions of P\'olya's theorem. Version arXiv:2511.07342v1 (CC BY 4.0), distributed under the Creative Commons Attribution 4.0 International licence, as recorded in the arXiv licence metadata for this exact version and shown on the abstract page https://arxiv.org/abs/2511.07342v1. Full rights notice: licenses/arxiv-2511.07342-CC-BY-4.0.txt.\par\smallskip

\noindent\textit{Editorial scope.} Counted unit: the Proposition whose source label is \texttt{prop:interiorCAcox} in the article (source0.tex lines 663--671, source label \texttt{prop:interiorCAcox}), delivered together with the article's own complete original proof of it (source0.tex lines 674--714, one \texttt{proof} environment of 41 lines). No mathematics is written, repaired or re-derived: statement and proof are the article's own text line for line. Required context retained uncounted: Prop:InteriorOfCA\_bis (lines 363--372); Eq:CoxHom (lines 542--545); Lemma:CoxSparseCopisitive (lines 557--560); Eq:JFaces (lines 605--608); lem:coxtruncation (lines 611--620); eq:decompositionIrrelevantLocus (lines 645--648); Lemma:VanishingOnIrrelevant (lines 650--653). Every definition, display and label of those lines is retained unchanged. Omitted from this package and not redistributed: 1--362, 373--541, 546--556, 561--604, 609--610, 621--644, 649--649, 654--662, 672--673, 715--1565. Nothing inside a retained excerpt is omitted or altered: every retained body is delivered exactly as the article's lines are written, so no omission receipt of any kind accompanies this package. Citations: the retained excerpts carry no citation command at all, so no bibliography entry of the article is transcribed and no reference is redistributed. Reference adaptation: none was needed and none was made; every reference inside the delivered unit resolves inside it, so no unresolved reference or citation is produced. Printed slips kept literal: no printed word, symbol, hypothesis or number of the counted statement or of its proof was altered. Layout and class: the article's own class is \texttt{amsart} with packages including geometry, inputenc, babel, amssymb, mathtools, amsfonts, amsmath, amsthm, dsfont, soul, biblatex, hyperref, tikz, tikz-cd; the wrapper is the lane's pinned \texttt{amsart} editorial wrapper at 10pt on A4, which restates the theorem-like environments the retained text needs and adds the article's own macro definitions that the retained text needs (\texttt{\textbackslash RR}, \texttt{\textbackslash conv}, \texttt{\textbackslash cox}), with extra wrapper packages (bm, bbm, dsfont, xspace, cancel, mathtools). The delivered numbering is the unit's own. Assets: no retained span contains a figure environment, a graphics-inclusion command, a PDF-page inclusion, a table or a TikZ picture; no image file is copied into this package, the package declares no asset dependency and no image file is redistributed with it. Licence: version arXiv:2511.07342v1 of this article is distributed under the Creative Commons Attribution 4.0 International licence, as recorded in the arXiv licence metadata for this exact version and shown on the abstract page https://arxiv.org/abs/2511.07342v1. A full rights notice is delivered at licenses/arxiv-2511.07342-CC-BY-4.0.txt. Characters: every retained excerpt is pure ASCII, so the delivered engine, XeLaTeX with its default Latin Modern text font, reports no missing character.
\begin{proposition}
\label{Prop:InteriorOfCA_bis}
Let $A \subseteq \mathbb{Z}^n$ be a finite set, $k \in \mathbb{N}$ and $f \in \RR[t_1^{\pm 1}, \dots , t_n^{\pm 1}]_{k \cdot A}$. 
Then the following are equivalent.
\begin{itemize}
    \item[(ii)] $f$ is strictly  $A$-copositive,
    \item[(iii)]  $f^{k\cdot(F \cap A)}(t) > 0$ for all $t \in \mathbb{R}^n_{>0}$ and all nonempty faces $F \subseteq \conv(A)$,
    \item[(iv)]  $f^G(t) > 0$ for all $t \in \mathbb{R}^n_{>0}$ and all nonempty faces $G \subseteq \conv(k \cdot A)$.
\end{itemize}
\end{proposition}

\begin{align}
\label{Eq:CoxHom}
    f_\cox(x)\coloneqq x^b f\big(\varphi_{F^\top}(x)\big) = \sum_{a\in A}c_a x^{F^\top a+b}\in \RR[x_1,\dots,x_r],
\end{align}

\begin{lemma}
\label{Lemma:CoxSparseCopisitive}
Let $A \subseteq \mathbb{Z}^n$ be a finite set such that $\conv(A)$ has dimension $n$. For $f \in \RR[t_1^\pm , \dots, t_n^\pm]_A$ denote $f_{\cox}$ its Cox homogenization as in \eqref{Eq:CoxHom}. Then $f$ is $A$-copositive if and only if $f_\cox(x) \geq 0$ for all $x \in \mathbb{R}^r_{\geq0}$. 
\end{lemma}

\begin{align}
\label{Eq:JFaces}
 G_I \coloneqq \{ \, a \in \RR^n \, \mid \, F_i \cdot a + b_i = 0 \text{ for all } i \in I \text{ and }  F_i \cdot a + b_i > 0 \text{ for all } i \in [r] \setminus I \, \}
 \end{align}

\begin{lemma}
\label{lem:coxtruncation}
Let $A \subseteq \mathbb{Z}^n$ be a finite set such that $\conv(A)$ has dimension $n$.
Let $I \subseteq [r]$ be the set of active constraints for the face $G_I$ of $\conv(A)$ as in \eqref{Eq:JFaces}.
For any $x \in \RR^r_{>0}$ we write $x_I \in \mathbb{R}^r_{\geq0}$ for the vector whose entries $(x_I)_i = 0$ if $i \in I$ and keep  $(x_I)_i = x_i$ if $i \in [r] \setminus I$.
For $f \in \RR[t_1^\pm , \dots, t_n^\pm]_A$ and $x\in \RR^r_{>0}$, we have
\begin{align*}
        f_{\cox}(x_I)=x^bf^{G_I}(\varphi_{F^\top}(x)).
    \end{align*}
\end{lemma}

\begin{align}
\label{eq:decompositionIrrelevantLocus}
    Z_A=\bigcup_{j=1}^k\mathcal{V}(x_i\, \mid \, i\in C_j).
\end{align}

\begin{lemma}
\label{Lemma:VanishingOnIrrelevant}
For $f \in \RR[t_1^\pm , \dots, t_n^\pm]_A$, we have that $f_\cox(x)=0$ for all points in the vanishing locus of the irrelevant ideal $x \in Z_A$. 
\end{lemma}

\begin{proposition}
\label{prop:interiorCAcox}
Let $A \subseteq \mathbb{Z}^n$ be a finite set such that $\conv(A)$ has dimension $n$. 
  A polynomial $f \in \RR[t_1^\pm , \dots, t_n^\pm]_A$ is strictly $A$-copositive if and only if 
  both of the following two conditions hold: \begin{enumerate}
        \item[(i)] $f_\cox$ is copositive and
        \item[(ii)] $\mathcal{V}(f_\cox) \cap \RR^r_{\ge 0} = Z_A \cap \RR^r_{\ge 0}$.
    \end{enumerate}
\end{proposition}

\begin{proof}
By Lemma~\ref{Lemma:CoxSparseCopisitive}, we have that $f$ is $A$-copositive if and only if condition (i) holds. 
Thus, it is enough to prove that an $A$-copositive polynomial $f$ is strictly $A$-copositive if and only if condition (ii) holds. By Proposition \ref{Prop:InteriorOfCA_bis}, we have that $f$ is strictly $A$-copositive if and only if 
\begin{align}
\label{eq:recallTruncationsforCox}
f^{G}(t)>0\text{ for all nonempty faces }G \subseteq \conv(A) \text{ and all }t\in\RR_{>0}^n.
    \end{align}


In the remainder of the proof, we assume that $f$ is $A$-copositive.
To prove that (ii) implies strict $A$-copositivity, let $t \in \RR^n_{>0}$ and let $I \subseteq [r]$ be the set of active constraints for a nonempty face $G_I$ of $\conv(A)$  (cf.~\eqref{Eq:JFaces}).
Since $\varphi_{F^\top}$ maps $\RR_{>0}^r$ surjectively onto $\RR_{>0}^n$ (see the proof of Lemma~\ref{Lemma:CoxSparseCopisitive}), there exists $x\in \RR_{>0}^r$ such that $\varphi_{F^\top}(x)=t$.
By Lemma~\ref{lem:coxtruncation}, we have
\begin{align}
\label{eq:ProofProp4_3}
          f^{G_I}(t)=f^{G_I}(\varphi_{F^\top}(x))=x^{-b}f_{\cox}(x_I),
        \end{align}
        where $x_I$ is given by
        $(x_I)_i = 0$ if $i\in I$  and $(x_I)_i = x_i$ if $i \in [r] \setminus I$.
If $x_I \notin Z_A$, then \eqref{eq:ProofProp4_3}, together with conditions (i) and (ii), implies that \eqref{eq:recallTruncationsforCox} holds, which in turn implies that $f$ is strictly $A$-copositive.

To show that $x_I\notin Z_A$, it suffices to find one polynomial in the irrelevant ideal $B(A)$ which does not vanish at $x_I$. Pick any vertex $v\in G_I$ (here we use that $G_I$ is not the empty face) and let $\sigma \in \Sigma_A(n)$ be the corresponding maximal cone.
By construction, the rays $F_i$ are contained in the cone $\sigma$ for all $i \in I$.
Since $(x_I)_i = x_i > 0$ for $i  \notin \sigma$ by the choice of $x$, we have
\begin{align*}
        \prod_{i\notin \sigma}(x_I)_i=\prod_{i\notin \sigma}x_i>0.
    \end{align*}
    In particular, $x_I\notin Z_A$. This concludes the proof that (ii) implies strict $A$-copositivity.
    
     To prove the other implication, assume that $f$ is strictly $A$-copositive or equivalently that \eqref{eq:recallTruncationsforCox} holds.
     By Lemma~\ref{Lemma:VanishingOnIrrelevant}, we have that $f_\cox(x) = 0$ for all $x \in Z_A$.
     Thus, we only have to show that $f_\cox(z)=0$ for some $z\in \RR_{\geq 0}^r$ implies $z\in Z_A$. 

Let $z \in \RR_{\geq 0}^r$ such that $f_\cox(z) = 0$. Using the notation $I = \{ i \in [r] \mid z_i = 0\} \subseteq [r]$, we associate to $z$ a point $x \in \RR^r_{>0}$, which is defined as $x_i = z_i$ for $i \in [r] \setminus I$ and $x_i = 1$ for $i \in I$. With the notation used in the first part of the proof, we have $z = x_I$.

If the facet normal vectors $\{F_i \mid i \in I\}$ formed a cone in $\Sigma_A$, then this cone would be dual to a face $G_I \subseteq\conv(A)$, and Lemma \ref{lem:coxtruncation} would imply that \begin{align*}
        0=f_\cox(x_I) =x^b f_{|G_I}(\varphi_{F^\top}(x)),
    \end{align*}
    which would contradict \eqref{eq:recallTruncationsforCox}.
    Thus, the normal vectors $\{F_i \mid i \in I\}$ don't span a cone in $\Sigma_A$ and therefore $I$ must contain some primitive collection $C_j \subseteq [r]$. Using \eqref{eq:decompositionIrrelevantLocus}, we have that $x_I \in \mathcal{V}(x_i \mid i\in C_j)\subseteq Z_A$ finishing the proof. 
\end{proof}

\end{document}
